\begin{abstract}
Optical integrated sensing and communication (O-ISAC) is emerging across
fiber, free-space optical, visible-light, LiFi, and photonic-terahertz systems,
yet differences in physical platforms, metric definitions, measurement
planes, and validation practices impede a unified assessment. This study
presents a PRISMA 2020-grounded, cross-modality systematic review and
metric-governed narrative synthesis of O-ISAC research published from January
2020 through 22 June 2026. From 1,733 identified records, 272 full-text reports
were assessed; 227 eligible reports were mapped to 206 unique studies. The
claim-governed evidence model contains 8,306 extracted claims, of which 8,203
support primary synthesis after excluding 31 context-only and 72 quarantined
claims. Photonic-terahertz (69 studies) and fiber (56) were the largest
modalities, followed by VLC/LiFi (38), FSO (31), hybrid optical (9), and other
optical systems (3). Resource sharing, shared hardware, and shared waveforms
were the dominant integration mechanisms. Across 404 trade-off records,
bandwidth/resource allocation and power/dynamic-range constraints were more
frequent than a narrow rate--resolution framing. Although 156 studies reached
at least laboratory or proof-of-concept validation, only 12 reported field or
deployment-oriented evidence; open data were available for 13 studies, and
code or models were unavailable or not reported for 197. A deterministic
technical appraisal rated 125 studies as strong overall contributors, but
benchmark readiness and direct comparison admissibility remained limited.
These findings position O-ISAC as a substantive 6G research domain while
identifying measurement-plane contracts, open artifacts, common baselines,
and field validation as the principal requirements for reproducible
cross-modality progress.
\end{abstract}

\begin{IEEEkeywords}
6G, benchmark readiness, fiber sensing, free-space optical communication,
integrated sensing and communication, LiFi, metric comparability, O-ISAC,
photonic terahertz, PRISMA, visible light communication.
\end{IEEEkeywords}
